\chapter{Diagrama de Secuencia}
\label{ch:secuencia}

Los diagramas de secuencia muestran la interacción entre el cliente Angular, la API Express, MongoDB y GridFS para los procesos principales. Las fuentes editables están en \path{docs/diagramas/secuencia}.

\section{P01: Registro e inicio de sesión}

\begin{figure}[H]
  \centering
  \includegraphics[width=\textwidth,height=0.85\textheight,keepaspectratio]{build/secuencia/secuencia-registro-e-inicio-sesion.pdf}
  \caption{Secuencia de registro, inicio de sesión y cierre local.}
  \label{fig:seq-auth}
\end{figure}

\section{P02 y P03: Crear receta y gestionar imagen}

\begin{figure}[H]
  \centering
  \includegraphics[width=\textwidth,height=0.85\textheight,keepaspectratio]{build/secuencia/secuencia-crear-receta-y-gestionar-imagen.pdf}
  \caption{Secuencia de creación de receta y gestión de su imagen.}
  \label{fig:seq-receta-imagen}
\end{figure}

\section{P05: Comentar y reseñar}

\begin{figure}[H]
  \centering
  \includegraphics[width=\textwidth,height=0.85\textheight,keepaspectratio]{build/secuencia/secuencia-comentar-y-resenar.pdf}
  \caption{Secuencia de comentarios y reseñas.}
  \label{fig:seq-comentarios}
\end{figure}

\section{P06: Seguir usuarios y consultar el feed}

\begin{figure}[H]
  \centering
  \includegraphics[width=\textwidth,height=0.85\textheight,keepaspectratio]{build/secuencia/secuencia-seguir-usuarios-y-feed.pdf}
  \caption{Secuencia de seguimiento, feed y campana provisional.}
  \label{fig:seq-feed}
\end{figure}
